
%%%===========================================
\section{Timeline}

\subsection{Research Schedule and Timeline}

According to the official academic calendar of the doctoral program, the student’s \textbf{enrollment began on January 25, 2024}, which also marks the starting point for the time count toward graduation requirements. Thus, \textbf{the maximum allowed period for submitting the final dissertation is October 25, 2028}, with the \textbf{deadline for registering for the qualifying examination set for June 1, 2026}. Within this timeframe, the research plan is structured to ensure consistent progress and early completion of key milestones, allowing time for feedback, adjustments, and deepening of the theoretical and empirical components as necessary.

\begin{table}[H]
\centering
\label{tab:timeline}
\renewcommand{\arraystretch}{1.3}
\footnotesize
\begin{tabularx}{\textwidth}{|c|X|}
\hline
\textbf{Period} & \textbf{Milestones and Planned Activities} \\
\hline
\textbf{2nd Semester 2025} & * Start of research activities. * Final refinement of the research design and hypotheses. * Initial literature review and consolidation of the theoretical basis for Chapter 1. * Drafting of the introduction and early theoretical articulation. \\
\hline
\textbf{Year 2026} & * Completion of Chapter 1. * Development of formal theoretical and mathematical models (Chapter 2). * Preparation for the qualifying exam (by June 2026). * Initial design of data structure and econometric strategy. \\
\hline
\textbf{Year 2027} & * Empirical validation and econometric modeling (Chapter 3). * Data collection and processing. * Testing of the Beta Models. * Integration of theoretical and empirical findings. \\
\hline
\textbf{Year 2028 (until Oct)} & * Final editing, adjustments, formatting, and submission of the dissertation. * Preparation for defense. * Optional deepening of case studies or production of related academic papers. \\
\hline
\end{tabularx}
\end{table}

This schedule was designed not only to meet the program's formal deadlines, but also to anticipate potential difficulties in modeling, data access, or empirical validation. The timeline includes deliberate margins to accommodate unforeseen circumstances and provides flexibility for refining the research, deepening promising analytical aspects, and exploring new hypotheses that may emerge throughout the doctoral process.


\newpage

% fim %